\chapter{Números reales y funciones}
\label{chap:numeros-reales-funciones}

El cálculo diferencial estudia cómo varían las cantidades. Antes de hablar de
límites o derivadas necesitamos precisar dos piezas del lenguaje: los números
con los que medimos y las funciones que relacionan unas cantidades con otras.
Este capítulo fija esa base con el detalle necesario para el curso, sin
construir los números reales ni repetir un curso completo de álgebra elemental.

\section{Sistemas numéricos y necesidad de los números reales}
\label{sec:sistemas-numericos}
\CanonicalUnit{CM-C01-U001}{}{}

Usaremos la convención
\[
  \N=\{1,2,3,\ldots\}, \qquad
  \Nzero=\{0,1,2,3,\ldots\}.
\]
Los enteros son
\(\Z=\{\ldots,-2,-1,0,1,2,\ldots\}\), y los racionales son los cocientes
\[
  \Q=\left\{\frac{p}{q}:p,q\in\Z,\ q\ne0\right\}.
\]
Cada ampliación responde a una operación que no siempre puede efectuarse en el
sistema anterior: restar dentro de \(\N\), por ejemplo, conduce a \(\Z\), y
dividir enteros conduce a \(\Q\) cuando el divisor no es cero.

Sin embargo, los racionales no bastan para describir todas las longitudes. La
diagonal de un cuadrado de lado uno tiene longitud \(\sqrt2\), y el resultado
siguiente muestra que esa longitud no pertenece a \(\Q\). Los números reales
\(\R\) completan la recta numérica: contienen a \(\Q\), incorporan números
irracionales como \(\sqrt2\) y no presentan los huecos de \(\Q\) que impedirían
formular con seguridad muchos argumentos de existencia.

La cadena
\[
  \N\subset\Z\subset\Q\subset\R
\]
es una orientación útil, no una invitación a estudiar aquí el tamaño o la
numerabilidad de estos conjuntos. Esos asuntos no son necesarios para los
objetivos de cálculo de este capítulo.

\subsection{Un primer modelo de demostración por contradicción}
\CanonicalUnit{CM-C01-U002}{}{}
\ProofRole{proof-002}{}

Una \emph{demostración por contradicción} comienza suponiendo que la afirmación
que queremos probar es falsa. Si esa suposición obliga a aceptar una
imposibilidad lógica, concluimos que la afirmación original era verdadera.

\begin{teorema}[Irracionalidad de \(\sqrt2\)]
  El número \(\sqrt2\) no es racional.
\end{teorema}

\begin{proof}
Supongamos, en contra de lo que queremos demostrar, que
\(\sqrt2=m/n\) para ciertos enteros \(m,n\), con \(n\ne0\). Podemos elegir la
fracción reducida, de modo que \(m\) y \(n\) no tengan ningún factor común
mayor que uno. Al elevar al cuadrado obtenemos
\[
  m^2=2n^2.
\]
Por tanto, \(m^2\) es par. El cuadrado de un entero impar es impar, así que
\(m\) debe ser par; escribimos \(m=2k\). Sustituyendo,
\[
  4k^2=2n^2, \qquad\text{y por tanto}\qquad n^2=2k^2.
\]
El mismo razonamiento muestra que \(n\) también es par. Entonces \(m\) y \(n\)
son divisibles por dos, lo que contradice que \(m/n\) estuviera reducida. La
suposición inicial es imposible y, en consecuencia, \(\sqrt2\notin\Q\).
\end{proof}

La estructura del argumento importa tanto como el ejemplo: se declara la
suposición contraria, se encadenan consecuencias justificadas y se identifica
con precisión la contradicción final.

\section{La recta real: orden, completitud y densidad}
\label{sec:recta-real}
\CanonicalUnit{CM-C01-U003}{}{}
\ProofRole{proof-001}{}

En este curso trabajaremos con \(\R\) como un \emph{cuerpo ordenado completo}.
La expresión ``cuerpo ordenado'' resume que podemos sumar, restar, multiplicar
y dividir por números no nulos, y que esas operaciones son compatibles con el
orden. Por ejemplo, si \(x<y\), entonces \(x+z<y+z\); si además \(z>0\),
entonces \(xz<yz\).

La palabra \emph{completo} expresa una propiedad que distingue a \(\R\) de
\(\Q\). La formularemos después de introducir supremos, pero su idea puede
anticiparse así: un conjunto no vacío de reales que no puede crecer sin límite
tiene una frontera superior real bien definida. Adoptamos esta propiedad como
fundamento de trabajo; no construiremos \(\R\) a partir de \(\Q\).

Entre dos números reales distintos siempre hay tanto un racional como un
irracional. Se dice que \(\Q\) y \(\R\setminus\Q\) son \emph{densos} en
\(\R\). Esta afirmación explica por qué racionales e irracionales aparecen
entremezclados en la recta, pero no elimina su diferencia: por ejemplo,
\(\sqrt2\) sigue siendo irracional aunque pueda aproximarse por racionales con
tanta precisión como se desee.

\section{Intervalos, distancia y entornos}
\label{sec:intervalos-distancia-entornos}
\CanonicalUnit{CM-C01-U004}{}{}

Un \emph{intervalo} es un conjunto de reales que contiene todos los puntos
situados entre cualesquiera dos de sus elementos. Para \(a<b\), escribimos
\[
\begin{aligned}
 (a,b)&=\{x\in\R:a<x<b\},
& [a,b]&=\{x\in\R:a\le x\le b\},\\
 (a,b]&=\{x\in\R:a<x\le b\},
& [a,b)&=\{x\in\R:a\le x<b\}.
\end{aligned}
\]
Los intervalos no acotados se escriben, por ejemplo,
\((a,+\infty)\), \(( -\infty,b]\) y \(\R=(-\infty,+\infty)\). Los símbolos
\(\pm\infty\) no representan números reales y por ello nunca se incluyen con
un corchete.

\begin{definicion}[Valor absoluto y distancia]
Para \(x\in\R\), el \emph{valor absoluto} de \(x\) es
\[
 |x|=\begin{cases}x,&x\ge0,\\-x,&x<0.\end{cases}
\]
La distancia entre \(x\) e \(y\) es \(d(x,y)=|x-y|\).
\end{definicion}

Así, \(|x-a|<r\), con \(r>0\), significa que la distancia de \(x\) a \(a\)
es menor que \(r\), y equivale a
\[
  a-r<x<a+r.
\]
Esta equivalencia conecta el lenguaje algebraico del valor absoluto con la
geometría de la recta.

\begin{definicion}[Entorno]
Si \(a\in\R\) y \(r>0\), el \emph{entorno abierto} de centro \(a\) y radio
\(r\) es
\[
  V_r(a)=(a-r,a+r)=\{x\in\R:|x-a|<r\}.
\]
El \emph{entorno reducido} es
\(V_r(a)\setminus\{a\}\).
\end{definicion}

Un entorno es, por tanto, un tipo particular de intervalo abierto con un centro
y un radio declarados. No usaremos las palabras \emph{intervalo} y
\emph{entorno} como si fueran intercambiables.

\ProofRole{proof-003}{}
\begin{proposicion}[Desigualdad triangular]
Para cualesquiera \(x,y\in\R\),
\[
  |x+y|\le |x|+|y|.
\]
\end{proposicion}

\begin{proof}
De \(-|x|\le x\le |x|\) y \(-|y|\le y\le |y|\) se obtiene, sumando,
\[
  -(|x|+|y|)\le x+y\le |x|+|y|.
\]
Para todo \(u\in\R\) y todo \(M\ge0\), la doble desigualdad
\(-M\le u\le M\) equivale a \(|u|\le M\). Aplicándola a
\(u=x+y\) y \(M=|x|+|y|\) resulta la afirmación.
\end{proof}

En adelante utilizaremos la desigualdad triangular como herramienta. Una forma
frecuente, obtenida aplicándola a \(x-y=(x-z)+(z-y)\), es
\[
  |x-y|\le |x-z|+|z-y|.
\]

\section{Cotas, extremos, supremo e ínfimo}
\label{sec:cotas-extremos}
\CanonicalUnit{CM-C01-U005}{}{}

\begin{definicion}[Cotas]
Sea \(S\subseteq\R\). Un número \(M\in\R\) es una \emph{cota superior} de
\(S\) si \(x\le M\) para todo \(x\in S\). Un número \(m\in\R\) es una
\emph{cota inferior} si \(m\le x\) para todo \(x\in S\). El conjunto es
\emph{acotado superiormente}, \emph{acotado inferiormente} o
\emph{acotado} según tenga cotas del tipo correspondiente o de ambos tipos.
\end{definicion}

Una cota no tiene por qué pertenecer al conjunto y rara vez es única. Por
ejemplo, \(1,2\) y \(100\) son cotas superiores de \((0,1)\), pero ninguno de
esos hechos afirma todavía cuál es la menor cota superior.

\begin{definicion}[Máximo y mínimo]
Un elemento \(M\in S\) es el \emph{máximo} de \(S\) si \(x\le M\) para todo
\(x\in S\). Análogamente, \(m\in S\) es el \emph{mínimo} si \(m\le x\) para
todo \(x\in S\). Si existen, son únicos y se denotan \(\max S\) y \(\min S\).
\end{definicion}

La pertenencia es decisiva: \([0,1]\) tiene máximo \(1\), mientras que
\([0,1)\) no tiene máximo. En efecto, dado \(x<1\), el número
\((x+1)/2\) sigue perteneciendo a \([0,1)\) y es mayor que \(x\).

\begin{definicion}[Supremo e ínfimo]
Sea \(S\subseteq\R\) no vacío. Si \(S\) está acotado superiormente, su
\emph{supremo}, denotado \(\sup S\), es la menor de sus cotas superiores. Si
\(S\) está acotado inferiormente, su \emph{ínfimo}, denotado \(\inf S\), es la
mayor de sus cotas inferiores.
\end{definicion}

El supremo puede pertenecer o no al conjunto. Para \(S=[0,1)\),
\(\sup S=1\), pero \(S\) no tiene máximo. En cambio, si \(S=[0,1]\), entonces
\(\sup S=\max S=1\). Del mismo modo, \(\inf(0,1]=0\), aunque ese intervalo no
tenga mínimo.

Una caracterización práctica es la siguiente: \(s=\sup S\) si y solo si
\(x\le s\) para todo \(x\in S\) y, para cada \(\varepsilon>0\), existe
\(x_\varepsilon\in S\) tal que \(s-\varepsilon<x_\varepsilon\). La segunda
condición dice que ninguna cantidad positiva puede rebajar \(s\) y dejarlo como
cota superior.

\ProofRole{proof-001}{}
\begin{teorema}[Propiedad del supremo o completitud de \(\R\)]
Todo subconjunto no vacío de \(\R\) que esté acotado superiormente tiene
supremo en \(\R\).
\end{teorema}

Esta propiedad se adopta como resultado de referencia. Su versión inferior se
deduce aplicándola al conjunto \(-S=\{-x:x\in S\}\): todo conjunto no vacío y
acotado inferiormente tiene ínfimo. Más adelante, la completitud sostendrá
resultados de existencia; no debe confundirse con una regla para calcular
automáticamente extremos.

\section{Funciones: objetos, no solo fórmulas}
\label{sec:funciones-objetos}
\CanonicalUnit{CM-C01-U006}{}{}

\begin{definicion}[Función]
Una \emph{función} \(f:A\to B\) asigna a cada elemento \(x\in A\) un único
elemento \(f(x)\in B\). El conjunto \(A\) es el \emph{dominio} y \(B\) es el
\emph{codominio}.
\end{definicion}

La función incluye la regla de asignación, el dominio y el codominio. Por eso
una fórmula aislada no determina por sí sola una función. Por ejemplo,
\[
 f:\R\to\R,\quad f(x)=x^2,
 \qquad\text{y}\qquad
 h:[0,\infty)\to[0,\infty),\quad h(x)=x^2
\]
usan la misma expresión, pero son funciones distintas. La primera no es
inyectiva ni sobreyectiva; la segunda es biyectiva.

En una fórmula como \(r(x)=1/(x-2)\), si no se declara otro dominio,
entenderemos el \emph{dominio natural} como el mayor subconjunto de \(\R\) en
el que la expresión está definida: aquí, \(\R\setminus\{2\}\). En un modelo,
en cambio, el dominio puede ser menor por razones vinculadas al significado de
la variable.

\subsection{Imagen, preimagen y gráfica}
\CanonicalUnit{CM-C01-U007}{}{}

\begin{definicion}[Imagen]
Si \(f:A\to B\) y \(E\subseteq A\), la \emph{imagen} de \(E\) es
\[
  f(E)=\{f(x):x\in E\}\subseteq B.
\]
En particular, \(f(A)\) es la imagen de la función. También puede llamarse
\emph{recorrido}, con este significado y nunca como sinónimo de codominio.
\end{definicion}

Siempre se cumple \(f(A)\subseteq B\), pero no necesariamente
\(f(A)=B\). Para \(f:\R\to\R\), \(f(x)=x^2\), el codominio es \(\R\) y la
imagen es \([0,\infty)\).

\begin{definicion}[Preimagen]
Si \(F\subseteq B\), la \emph{preimagen} de \(F\) es
\[
  f^{-1}(F)=\{x\in A:f(x)\in F\}.
\]
Esta notación está definida para cualquier función; no presupone que exista una
función inversa.
\end{definicion}

Para la función cuadrado \(f:\R\to\R\), por ejemplo,
\[
  f^{-1}([1,4])=[-2,-1]\cup[1,2].
\]
Esta preimagen es un conjunto, no el valor de una función inversa.

La \emph{gráfica} de \(f:A\to B\) es
\[
  \operatorname{Gr}(f)=\{(x,f(x)):x\in A\}\subseteq A\times B.
\]
La gráfica representa una función ya definida; no sustituye la información
sobre su dominio y su codominio.

\subsection{Inyectividad, sobreyectividad e inversa}
\CanonicalUnit{CM-C01-U008}{}{}

Una función \(f:A\to B\) es
\begin{itemize}
  \item \emph{inyectiva} si \(f(x_1)=f(x_2)\) implica \(x_1=x_2\);
  \item \emph{sobreyectiva} si \(f(A)=B\);
  \item \emph{biyectiva} si es inyectiva y sobreyectiva.
\end{itemize}

\begin{proposicion}[Función inversa]
Una función \(f:A\to B\) posee una función inversa
\(f^{-1}:B\to A\) si y solo si es biyectiva. En ese caso,
\[
  f^{-1}\circ f=\id_A,
  \qquad
  f\circ f^{-1}=\id_B.
\]
\end{proposicion}

La invertibilidad depende de la función completa, no solo de la expresión. La
función \(x\mapsto x^2\) no tiene inversa como función de \(\R\) en \(\R\),
pero su restricción \(h:[0,\infty)\to[0,\infty)\) sí la tiene y
\(h^{-1}(y)=\sqrt y\).

\subsection{Composición correctamente tipada}
\CanonicalUnit{CM-C01-U009}{}{}

Sean
\[
  f:A\to B, \qquad g:C\to D.
\]
Para aplicar primero \(f\) y después \(g\), cada salida efectiva de \(f\) debe
ser una entrada admisible de \(g\). Por tanto, la condición necesaria y
suficiente para definir la composición en todo \(A\) es
\[
  f(A)\subseteq C.
\]
Bajo esta compatibilidad se define la \emph{composición}
\[
  g\circ f:A\to D, \qquad (g\circ f)(x)=g(f(x)).
\]
La inclusión relevante compara la imagen de \(f\) con el dominio de \(g\); no
compara codominios sin atender a los valores realmente producidos.

\begin{ejemplo}
Sean \(f:\R\to[0,\infty)\), \(f(x)=x^2\), y
\(g:[0,\infty)\to\R\), \(g(u)=\sqrt u\). Como
\(f(\R)=[0,\infty)\), la composición está definida y
\((g\circ f)(x)=|x|\). En cambio,
\((f\circ g)(u)=u\) para \(u\ge0\). El orden de composición importa.
\end{ejemplo}

Si una regla \(g\) solo está definida en una parte de \(f(A)\), puede
considerarse la composición en el subconjunto
\(\{x\in A:f(x)\in C\}\), pero entonces ese subconjunto, y no todo \(A\), es
su dominio. Esta observación evita ocultar restricciones de dominio.

\section{Revisión selectiva de funciones elementales}
\label{sec:funciones-elementales}

Las familias siguientes aparecerán repetidamente en cálculo. El objetivo aquí
es fijar dominios, imágenes, relaciones inversas y rasgos gráficos esenciales,
no reconstruir toda la teoría elemental.

\subsection{Polinómicas, racionales y definidas a trozos}
\CanonicalUnit{CM-C01-U010}{}{}

Una función polinómica tiene la forma
\[
  p(x)=a_0+a_1x+\cdots+a_nx^n,
\]
donde \(n\in\Nzero\), los coeficientes son reales y, si \(n\ge1\),
\(a_n\ne0\). Su dominio natural es \(\R\). Su comportamiento global depende
del grado y del signo del coeficiente principal, aunque el análisis fino se
hará más adelante con herramientas de cálculo.

Una función racional está dada por
\[
  r(x)=\frac{p(x)}{q(x)},
\]
donde \(p\) y \(q\) son polinomios y \(q\) no es el polinomio nulo. Su dominio
natural es
\[
  \{x\in\R:q(x)\ne0\}.
\]
No se puede cancelar un factor y olvidar los puntos excluidos por el
denominador original: \((x^2-1)/(x-1)\) coincide con \(x+1\) cuando
\(x\ne1\), pero la primera expresión no está definida en \(x=1\).

Una función definida a trozos utiliza reglas distintas en partes declaradas
del dominio. El valor absoluto es el ejemplo básico. Para que la definición
sea una función, cada punto del dominio debe recibir exactamente un valor; si
los trozos se solapan, sus reglas han de ser compatibles en el solapamiento.

\subsection{Exponenciales y logaritmos}
\CanonicalUnit{CM-C01-U011}{}{}
\ProofRole{proof-004}{}

Para una base \(a>0\), la función exponencial de base \(a\) es
\[
  \exp_a:\R\to(0,\infty), \qquad \exp_a(x)=a^x.
\]
La extensión de las potencias desde exponentes racionales hasta exponentes
reales requiere la completitud de \(\R\); adoptamos aquí esa construcción y
sus propiedades básicas como resultado de referencia. En particular,
\[
  a^{x+y}=a^x a^y, \qquad a^{-x}=\frac1{a^x}.
\]
Si \(a>1\), la exponencial es estrictamente creciente; si \(0<a<1\), es
estrictamente decreciente. En ambos casos es biyectiva de \(\R\) en
\((0,\infty)\). Para \(a=1\) es constante y no posee inversa.

El número \(e\) es la base distinguida de la exponencial natural. En este
capítulo se toma como una constante real positiva bien definida; sus
caracterizaciones mediante límites, derivadas o series se estudiarán solo
cuando esas herramientas estén disponibles.

Si \(a>0\) y \(a\ne1\), el \emph{logaritmo de base \(a\)} es la función
inversa de la exponencial:
\[
  \log_a:(0,\infty)\to\R,
  \qquad
  \log_a y=x \iff a^x=y.
\]
El \emph{logaritmo natural} se denota \(\ln=\log_e\). No utilizaremos el
símbolo \(\log\) sin declarar su base.

Las relaciones inversas son
\[
  \log_a(a^x)=x \quad(x\in\R),
  \qquad
  a^{\log_a y}=y \quad(y>0).
\]
Las restricciones escritas forman parte de las identidades.

\ProofRole{proof-005}{}
\begin{proposicion}[Identidades logarítmicas representativas]
Si \(a>0\), \(a\ne1\), y \(x,y>0\), entonces
\[
 \log_a(xy)=\log_a x+\log_a y,
 \qquad
 \log_a\!\left(\frac{x}{y}\right)=\log_a x-\log_a y.
\]
Además, para \(r\in\R\), \(\log_a(x^r)=r\log_a x\).
\end{proposicion}

\begin{proof}
Demostraremos la primera identidad como modelo. Sean
\(u=\log_a x\) y \(v=\log_a y\). Entonces \(x=a^u\), \(y=a^v\), y
\[
  xy=a^u a^v=a^{u+v}.
\]
Al aplicar \(\log_a\) a ambos lados, que son positivos, resulta
\(\log_a(xy)=u+v=\log_a x+\log_a y\). Las otras identidades se obtienen del
mismo modo usando las leyes de los exponentes y respetando las hipótesis de
positividad.
\end{proof}

\begin{advertencia}
Expresiones como \(\ln(x-y)\) o \(\log_a(q(x))\) solo tienen sentido real en
los puntos donde el argumento es positivo; además, una base logarítmica real
debe cumplir \(a>0\) y \(a\ne1\). Las identidades algebraicas no eliminan
estas condiciones.
\end{advertencia}

\subsection{Funciones trigonométricas y radianes}
\CanonicalUnit{CM-C01-U012}{}{}
\ProofRole{proof-005}{}

La medida en \emph{radianes} identifica un ángulo orientado con la longitud de
arco que determina en la circunferencia unidad. Esta elección hace que una
vuelta completa corresponda a \(2\pi\) y es la medida compatible con las
fórmulas fundamentales del cálculo.

Para cada \(t\in\R\), el punto de la circunferencia unidad alcanzado tras
recorrer un arco orientado de longitud \(t\) tiene coordenadas
\((\cos t,\sin t)\). Así se definen
\[
  \sin,\cos:\R\to[-1,1]
\]
para todo número real, no solo para ángulos de un triángulo. Ambas funciones
son periódicas de periodo \(2\pi\); el coseno es par y el seno es impar. La
identidad pitagórica
\[
  \sin^2 t+\cos^2 t=1
\]
se obtiene directamente de la ecuación \(x^2+y^2=1\) de la circunferencia.

Las funciones cociente tienen dominios restringidos:
\[
  \tan t=\frac{\sin t}{\cos t}
  \quad\text{si }\cos t\ne0,
  \qquad
  \cot t=\frac{\cos t}{\sin t}
  \quad\text{si }\sin t\ne0.
\]
En particular,
\[
 \Dom(\tan)=\R\setminus\left\{\frac\pi2+k\pi:k\in\Z\right\},
 \qquad
 \Dom(\cot)=\R\setminus\{k\pi:k\in\Z\}.
\]

Como las funciones trigonométricas son periódicas, no son inyectivas en todo
\(\R\). Para definir inversas se restringen a intervalos donde sí lo son:
\[
\begin{aligned}
 \arcsin&:[-1,1]\to[-\pi/2,\pi/2],\\
 \arccos&:[-1,1]\to[0,\pi],\\
 \arctan&:\R\to(-\pi/2,\pi/2).
\end{aligned}
\]
Por ejemplo, \(\sin(\arcsin y)=y\) para \(y\in[-1,1]\), mientras que
\(\arcsin(\sin t)=t\) solo cuando \(t\in[-\pi/2,\pi/2]\). Esta asimetría es
otra manifestación de que una relación inversa depende del dominio escogido.

Los gráficos cualitativos completan estas descripciones: las polinómicas están
definidas en todo \(\R\); las racionales pueden presentar puntos excluidos; las
exponenciales toman valores positivos; los logaritmos solo admiten entradas
positivas; y las trigonométricas repiten su patrón. El análisis mediante
límites, continuidad y derivadas pertenece a capítulos posteriores y no se
presupone aquí.
